% ==============================================================================
% SECCIÓN 3: REGLAS DE NEGOCIO DEL SISTEMA (RN-xx)
% ==============================================================================

\section{Reglas de Negocio del Sistema (RN)}

Las Reglas de Negocio definen las restricciones operacionales, financieras, legales y de seguridad que rigen el comportamiento de la empresa Comercial Machi's y que deben ser implementadas y validadas por el sistema SYSELL tanto en el frontend como a nivel de dominio en el backend.

\subsection{Especificación Formal de Reglas de Negocio}

\begin{rnbox}{RN-01}{Moneda Nacional (PEN) y Régimen Tributario IGV (18\%)}
\textbf{Descripción:} Toda transacción comercial, precio de venta al público, costo unitario, descuento, monto de apertura de caja y balance financiero debe ser registrado y procesado exclusivamente en \textbf{Soles Peruanos (PEN, \texttt{S/.})}.\\
\textbf{Fórmula de Cálculo Tributario:}
\begin{align}
    \text{Base Imponible} &= \frac{\text{Precio Total}}{1 + 0.18} = \frac{\text{Precio Total}}{1.18} \\
    \text{Monto IGV} &= \text{Precio Total} - \text{Base Imponible} = \text{Base Imponible} \times 0.18
\end{align}
\textbf{Restricción:} Los comprobantes de venta deben mostrar obligatoriamente el desglose de la Base Imponible, el IGV calculado (18\%) y el Importe Total con precisión de dos decimales (redondeo bancario al centésimo más cercano).
\end{rnbox}

\begin{rnbox}{RN-02}{Integridad de Stock, Transaccionalidad y Prohibición de Saldos Negativos}
\textbf{Descripción:} El sistema no permitirá bajo ninguna circunstancia el registro de una venta, salida o transferencia si la cantidad solicitada supera el stock físico disponible en el local de origen.\\
\textbf{Validación Concurrente:} Toda operación que afecte el inventario debe ejecutarse dentro de una transacción de base de datos (\texttt{DB::transaction}) con bloqueo pesimista a nivel de fila (\texttt{SELECT ... FOR UPDATE}), evitando condiciones de carrera (\textit{race conditions}) entre terminales concurrentes.
\end{rnbox}

\begin{rnbox}{RN-03}{Umbral de Stock Mínimo y Alertas Automáticas}
\textbf{Descripción:} Cada variante de producto (por combinación de modelo, talla, color y local) posee un campo parametrizable denominado \texttt{stock\_minimo} (valor entero $\ge 1$, por defecto 3 unidades).\\
\textbf{Disparador de Alerta:} Cuando el stock disponible en un local sea menor o igual al \texttt{stock\_minimo} ($\text{Stock} \le \text{Stock Mínimo}$), el sistema catalogará la variante en estado de \textbf{Alerta Amarilla / Roja} y generará una notificación visual inmediata en el panel del Administrador para ordenar la reposición.
\end{rnbox}

\begin{rnbox}{RN-04}{Inmutabilidad Transaccional y Borrado Lógico (Soft Delete)}
\textbf{Descripción:} Para preservar la trazabilidad fiscal, contable y operativa, queda estrictamente prohibida la eliminación física (\texttt{DELETE}) de registros en las tablas de: Usuarios, Productos, Variantes, Locales, Ventas, Detalles de Venta, Devoluciones y Cierres de Caja.\\
\textbf{Mecanismo:} Las bajas se implementarán mediante la columna \texttt{deleted\_at} (tipo \texttt{TIMESTAMP NULL}). Una entidad dada de baja se excluirá de las consultas de catálogo activo (\textit{scope query}), pero conservará todas sus relaciones históricas en reportes pasados.
\end{rnbox}

\begin{rnbox}{RN-05}{Restricción de Métodos de Pago Operativos}
\textbf{Descripción:} Las ventas únicamente pueden cancelarse mediante los tres métodos de pago autorizados en mostrador:
\begin{enumerate}[leftmargin=*,noitemsep]
    \item \textbf{Efectivo:} Requiere ingresar el \texttt{monto\_recibido} ($\ge \text{Total Venta}$). El sistema calcula automáticamente el vuelto: $\text{Vuelto} = \text{Monto Recibido} - \text{Total Venta}$.
    \item \textbf{Billetera Digital (Yape / Plin):} Requiere ingresar obligatoriamente el \texttt{codigo\_operacion} (de 6 a 8 dígitos numéricos) verificado en la pantalla del cliente.
    \item \textbf{Transferencia Interbancaria:} Requiere ingresar el \texttt{numero\_operacion} bancario y la entidad bancaria emisora.
\end{enumerate}
\end{rnbox}

\begin{rnbox}{RN-06}{Políticas de Devolución, Cambio de Prenda y Garantía}
\textbf{Descripción:} Un cliente puede solicitar cambio de talla/color o devolución de prenda bajo las siguientes condiciones:
\begin{enumerate}[leftmargin=*,noitemsep]
    \item Plazo máximo de 7 días calendario posteriores a la fecha de compra original.
    \item Presentación obligatoria del número de comprobante/nota de venta original.
    \item La prenda debe reintegrarse en estado apto al stock del local donde se realiza el cambio.
    \item Si el cambio es por un producto de mayor valor, el cliente abona la diferencia; si es de menor valor, se emite una \textbf{Nota de Crédito Interna} (no se efectúa devolución de efectivo en mostrador sin autorización explícita del Administrador).
\end{enumerate}
\end{rnbox}

\begin{rnbox}{RN-07}{Apertura, Arqueo Ciego y Cierre de Caja Diario por Local}
\textbf{Descripción:} 
\begin{enumerate}[leftmargin=*,noitemsep]
    \item Antes de registrar la primera venta del día, el cajero debe realizar la \textbf{Apertura de Caja} ingresando el \texttt{monto\_inicial} (fondo para vuelto).
    \item Al finalizar la jornada o turno, el cajero debe realizar el \textbf{Cierre de Caja} ingresando el arqueo físico de efectivo y el consolidado de comprobantes digitales (arqueo a ciegas).
    \item El sistema contrastará el conteo físico contra el cálculo del sistema ($\text{Esperado} = \text{Monto Inicial} + \sum \text{Ventas Efectivo} - \sum \text{Gastos Menores}$), calculando la diferencia (Sobrante, Faltante o Conforme) y bloqueando nuevas ventas en esa sesión de caja.
\end{enumerate}
\end{rnbox}

\begin{rnbox}{RN-08}{Seguridad de Acceso: Bloqueo por Intentos Fallidos y Caducidad de Sesión}
\textbf{Descripción:} 
\begin{enumerate}[leftmargin=*,noitemsep]
    \item Si un usuario introduce una contraseña incorrecta en 5 intentos consecutivos, su cuenta quedará bloqueada por un período de \textbf{15 minutos} para mitigar ataques de fuerza bruta.
    \item Si una sesión autenticada permanece inactiva por más de \textbf{15 minutos} (sin interacción en el navegador), el token JWT/Sanctum caduca automáticamente, redirigiendo al usuario al formulario de inicio de sesión.
\end{enumerate}
\end{rnbox}

\begin{rnbox}{RN-09}{Trazabilidad Universal y Registro Inmutable de Auditoría (Audit Logs)}
\textbf{Descripción:} Todo evento crítico del sistema (inicio de sesión, cambio de contraseña, creación de producto, modificación de precio, anulación de venta, baja lógica de usuario y cierre de caja) debe generar un registro automático e inmutable en la tabla \texttt{audit\_logs}, capturando: \texttt{user\_id}, \texttt{action}, \texttt{entity\_type}, \texttt{entity\_id}, \texttt{old\_values}, \texttt{new\_values}, \texttt{ip\_address} y \texttt{timestamp}.
\end{rnbox}

\begin{rnbox}{RN-10}{Variantes de Producto y Unicidad de Código SKU}
\textbf{Descripción:} Cada producto padre pertenece a una categoría y posee una marca y temporada. Cada combinación específica de (Producto, Talla, Color) genera una variante única con su propio código \textbf{SKU} (\textit{Stock Keeping Unit}) o código de barras irrepetible en todo el sistema.
\end{rnbox}
